\section{最小描述长度与 Bayesian 学习}
\label{sec:07-Information-Theory/06-Information-and-Learning/02}

十个数据点，九次多项式，训练残差严格为零。这个拟合无可挑剔，也毫无内容——它只是把那十个点换了一种写法。极端一点，最强的``模型''就是整张数据表本身：训练误差同样为零，泛化能力同样是零。

问题不在于拟合得太好，而在于我们一直只盯着账本的一栏。残差被压到零，模型那一栏于是膨胀了，可那一栏从来没被计价。最小描述长度（minimum description length，MDL）做的事就是把两栏放进同一个单位。

\subsection{把模型和残差记进同一本账}

设想你要把一份数据传给远方的同事，只按传输的比特数付费。你可以逐点直接发，也可以先发一个模型，再发``数据相对这个模型的偏离''。总账单是

\[
L(M)+L(D\given M).
\]

第一项是说清模型本身要花的比特，第二项是接收者拿到模型之后再描述数据要花的比特。模型太简单，第一项短而第二项长；模型太复杂，拟合项短了，模型说明却长得离谱。MDL 选的是让总长度最短的那个模型，也就是最能把数据里的规律榨成压缩量的模型。

这里有一处必须先堵住：码长不许事后编造。如果允许看完数据再免费设计编码，那就宣布当前这份数据的码字是空串，任何模型都能``完美压缩''。所以 MDL 要求编码预先约定且唯一可译，Kraft 不等式随即封死了同时给所有模型和所有数据都安排短码的可能，短码是稀缺资源，给了这边就没了那边。前缀码与 Kraft 不等式的完整交代见\hyperref[ch:057]{``无损信源编码：熵为什么是压缩极限''}。

有了这条约束，概率模型自动变成编码方案。分布 \(Q\) 诱导的理想码长是

\[
L(x)=-\log_2 Q(x),
\]

模型给数据的概率越高，码字越短。负对数似然从此不只是一个训练目标，它就是``在这个模型下描述这份数据''的长度。

\begin{center}
\begin{tikzpicture}[x=1.35cm,y=1.05cm,font=\small]
  \draw[->,black!55] (0,0) -- (6.3,0) node[right,font=\footnotesize] {模型复杂度};
  \draw[->,black!55] (0,0) -- (0,4.15) node[above,font=\footnotesize] {比特};
  \draw[blue!70,thick] plot[domain=0.15:5.7,samples=60] (\x,{0.45*\x});
  \draw[orange!85,thick] plot[domain=0.15:5.7,samples=90] (\x,{1.6/(\x+0.35)+0.4});
  \draw[black!85,very thick] plot[domain=0.15:5.7,samples=90]
    (\x,{0.45*\x+1.6/(\x+0.35)+0.4});
  \draw[black!30,dashed] (1.536,0) -- (1.536,1.94);
  \fill[red!75] (1.536,1.94) circle (1.6pt);
  \node[font=\footnotesize,black!75,right] at (4.35,1.72) {\(L(M)\)};
  \node[font=\footnotesize,black!75,right] at (4.35,0.62) {\(L(D\given M)\)};
  \node[font=\footnotesize,black!85,right] at (4.05,3.05) {总长度};
  \node[font=\footnotesize,black!70,below] at (1.536,-0.12) {最短处};
  \node[font=\footnotesize,black!55,align=center] at (1.15,3.55) {欠拟合：\\模型短，数据长};
  \node[font=\footnotesize,black!55,align=center] at (5.05,3.75) {过拟合：\\数据短，模型长};
\end{tikzpicture}

\emph{两条分量沿相反方向走，和有一个最小值。这张图承担了正则化的信息论读法：惩罚项不是从外面加进来的调味料，它就是账本上``说明模型''那一栏。最短处的位置由数据量决定——样本越多，橙色曲线整体抬高，最小值向右移，复杂模型才开始划算。}
\end{center}

\subsection{一枚硬币的十次投掷}

抛十次得到八次正面。比较两个模型：\(M_0\) 断言硬币严格公平，没有需要说明的参数；\(M_1\) 允许正面概率自由取值，用最大似然值 \(\hat p=0.8\) 拟合。

若编码的是具体出现顺序，\(M_0\) 给每一串的概率都是 \(2^{-10}\)，数据码长十比特。\(M_1\) 下的码长是

\[
-\log_2\bigl(0.8^{8}\cdot 0.2^{2}\bigr)\approx 7.22\ \text{比特}.
\]

灵活模型在拟合那一栏省下约 \(2.78\) 比特。可解码器还得知道用了哪个模型、参数是 \(0.8\)，以及这个 \(0.8\) 精确到第几位。这份说明若超过 \(2.78\) 比特，公平模型的总长度反而更短。样本继续增加而偏差持续存在时，拟合收益按样本量线性累积，模型说明的成本却基本不变，天平才会翻过去。

这个小例子顺带暴露一个常被跳过的麻烦：连续参数若要求无限精度，就得花无限比特。真做 MDL 必须指定量化精度、参数编码方式，或者干脆改用不依赖某个点估计的通用码。拿``参数个数''直接冒充模型码长只是一种粗糙近似，好用，但别当成定义。

\subsection{两段式 MDL 就是 MAP，于是正则化有了解释}

若约定

\[
L(M)=-\log_2 P(M),\qquad L(D\given M)=-\log_2 P(D\given M),
\]

总码长立刻变成 \(-\log_2\bigl[P(M)P(D\given M)\bigr]\)。负对数是单调递减的，最小化总长度等价于最大化 \(P(M)P(D\given M)\)，也就是最大后验选择，见\hyperref[sec:11-Mathematical-Statistics/02-Point-Estimation/04]{``MAP 把先验与似然合在一个点上''}。模型先验对应模型码长，似然对应给定模型后的数据码长——两套语言，一个式子。

这条对应把正则化从``经验做法''挪到了账本里。参数的 Gaussian 先验取负对数得到平方项，于是岭回归的 \(\lambda\norm{\theta}^2\) 就是说明这组参数所需的比特数；Laplace 先验取负对数得到绝对值项，Lasso 的稀疏偏好对应的是``大多数参数为零，因此不必逐个说明''这种编码。惩罚系数 \(\lambda\) 的位置也不再神秘：它是两栏之间的汇率，由先验的尖锐程度定下来。正则化路径与偏差方差的完整讨论在\hyperref[sec:11-Mathematical-Statistics/06-Resampling-and-Model-Selection/04]{``偏差---方差、正则化与模型选择''}。

不过等价止于此，不能顺势推出``MDL 就是 Bayesian 推断''。不同的 MDL 方案会采用归一化最大似然、预序列码或其他通用编码，压根不需要指定先验；Bayesian 那边则要求先验说清楚，而且关心的是完整后验而非一个 MAP 点。两者共用``概率即码长''这座桥，桥两头的目标并不相同，在有限样本上给出不同选择是常事。

\subsection{边际似然自带的 Occam 因子}

Bayesian 模型比较用的是证据

\[
P(D\given M)=\int P(D\given\theta,M)\,P(\theta\given M)\,d\theta,
\]

它是把参数的不确定性按先验平均之后的预测概率，\(-\log_2 P(D\given M)\) 可以读成一段混合码的长度。灵活模型也许在某个最佳参数处似然极高，但只要绝大多数参数区域预测得很差，积分就不会跟着高。

这就是 Occam 因子，而它的来源朴素得没有任何神秘感：概率守恒。一个模型若给许多互不相同的数据集都留了可能性，它就不可能同时给眼下这一份极高的概率——总量是一。所谓``简单的更可信''不需要额外的美学假设，它是归一化的后果。

同一件事换成比特来说：越灵活的模型，它的预测分布摊得越平，任何一份具体数据在它下面的码长就越长。

\subsection{AIC 与 BIC 读成两种账本}

模型选择准则常被当成公式背下来，放进描述长度的框架里则各有出处。

BIC 的形式 \(-2\log\hat L+k\log n\) 正是两段式 MDL 的渐近展开：\(k\log n\) 这一项来自把 \(k\) 个连续参数编码到 \(1/\sqrt{n}\) 量级的精度——数据越多，参数值得说得越细，模型那一栏就越贵。它的目标是在候选集合里选出真实模型，前提是真实模型确实在候选集合内。

AIC 的形式 \(-2\log\hat L+2k\) 走的是另一条路。它估计的是用拟合模型去预测新数据时的期望 KL 散度，惩罚项 \(2k\) 是拟合带来的乐观偏差的一阶修正。它不假设真实模型在候选集合里，追求的是预测风险最小。

两者常给出不同答案，这不是谁算错了。样本量大时 \(k\log n\) 远大于 \(2k\)，BIC 更保守；而当真实机制根本不在候选集合中时，AIC 那种``选一个预测得最准的近似''往往更合用。先说清目标是找真模型还是找好预测，再挑准则。

\subsection{短描述不自动等于好泛化}

可压缩的规律通常有复用价值，但 MDL 的结论依赖一串没写在公式里的约定：数据对象和顺序究竟怎样编码，模型类里有没有装进合理的机制，码表、参数精度与共享常数怎么算，样本是否独立同分布，未来的分布是否还和现在一样。任何一条改了，最短的那个模型都可能换人。

反例也不难想。一个短描述可能利用的是不稳定的采样偏差——训练集里恰好所有正例都来自同一个数据源，``按来源判''的编码极短，换个批次就崩。反过来，一个为因果解释设计的模型未必给出最短的预测码，因为它要付出说明机制的成本，而预测本身并不需要那份说明，这一点在\hyperref[sec:13-Machine-Learning/12-Causal-Inference/06]{``可识别不等于可外推也不等于可决策''}里有更细的展开。

所以 MDL 提供的是一个可以核对的复杂度账本，不是免于验证的泛化保证。它的实际价值在于：一旦你被迫写出 \(L(M)\)，那些平时藏在超参数里的假设就必须摆到台面上计价。

到这里，比特这把尺子量过了数据的价值和模型的成本。还剩一个位置没量：数据进模型之前先被压成的那个中间表示。下一节把``该保留输入的哪些部分''也写成两个互信息之间的取舍——顺带看看这个漂亮的目标函数在落地时遇到了什么麻烦。
